\subsection{Corrected ROS 2 Closed-Loop Verification Result}
\label{sec:ros-verification-results}

The corrected validation task ran on Ubuntu~24.04.4 with ROS~2 Jazzy,
Gazebo Harmonic, and RViz2. Gazebo supplied the station/chaser rendering and
camera image only; the ROS-native HCW node remained the state source of truth.
The input route was reconstructed from the corrected 24-node graph using the
complete glTF scene hierarchy and a conservative
$r_{\mathrm{veh}}=\unit{0.80}{m}$ chaser bounding radius. The route manifest
records 10 planned SOOAs, \unit{900}{s} mission duration, 81.33\% planned
coverage, and \unit{2.808}{m} minimum planned finite-body clearance above the
required \unit{2}{m} margin.

The 5$\times$ graphical closed-loop execution credited all 10 terminal
observations, reached 81.33\% coverage, and completed at mission time
\unit{900.01}{s} without a failed action. The post-run audit evaluated 18,918
recorded trajectory samples and every intervening segment against all 247,525
fully transformed ISS triangles. No segment crossed the mesh. The smallest
sampled center distance was \unit{8.506}{m}; after the continuous segment lower
bound and the \unit{0.80}{m} vehicle radius, the minimum body distance was
\unit{7.692}{m}, or \unit{5.692}{m} above the required safety margin. Peak
filtered acceleration was $0.0600~\mathrm{m/s^2}$. Terminal, reference-stream,
data-presence, acceleration, full-mesh, swept-crossing, and finite-body gates
all passed.

The retained evidence directory contains the six core CSV streams,
configuration and route manifests, launch log, summary, version-2 mesh audit,
raw RViz and camera streams, an annotated H.264 video, and a milestone preview.
The 198.3-s video uses RViz2 as the main view and the live
\texttt{/chaser/camera/image} topic as an inset. It is an execution and
interface demonstration, not an additional held-out comparison row.
Figure~\ref{fig:ros-key-camera-views} links all ten synchronized camera frames
to their exact credited positions in two orthogonal projections of the recorded
closed-loop trajectory.

\begin{figure}[t]
  \centering
  \includegraphics[width=0.95\linewidth]{figures/ros_key_camera_views_trajectory.pdf}
  \caption{Trajectory-linked camera views from the accepted corrected
  \texttt{validation\_002} task. Views 1--5 link locally to the radial--cross-track
  ($x$--$z$) projection, and views 6--10 link to the along-track--cross-track
  ($y$--$z$) projection; hollow gray markers retain the other observations as
  context. Inset labels give mission time and cumulative coverage. Dashed gray
  and solid blue denote planned and executed paths, and the star is the initial
  state. Apparent overlap in either projection does not imply contact: the
  annotated clearances and collision decision use the full 3-D transformed mesh
  and swept path. Gazebo rendered the views, while ROS HCW dynamics propagated
  state; coverage came from frozen geometric visibility masks, not image pixels.}
  \label{fig:ros-key-camera-views}
\end{figure}

The earlier paired ROS campaign used a translation-only interpretation of the
glTF node hierarchy. Re-auditing its illustrative route with the complete
hierarchy exposed a swept intersection with the S4 truss and a negative
finite-body clearance. Those numerical ROS campaign claims and their old video
are therefore excluded rather than carried into the corrected manuscript. The
current result establishes consistency of one complete closed-loop task with
the corrected geometry and audit; paired ROS method superiority would require a
new full campaign from the corrected route bundle.
